\begin{frame}{Positional Encoding}

  \begin{itemize}
    \item[$\blacktriangleright$] Transformers lack sequence order awareness.
    \item[$\blacktriangleright$] \textbf{Positional encodings} are added to input embeddings to provide information about token positions.
    \item[$\blacktriangleright$] \textbf{Sinusoidal encoding:}
  \end{itemize}
  
  \vspace{0.8em}
  
  \begin{align*}
    PE(pos, 2i) &= \sin\left(\frac{pos}{10000^{\frac{2i}{d_{\text{model}}}}}\right) \\
    PE(pos, 2i+1) &= \cos\left(\frac{pos}{10000^{\frac{2i}{d_{\text{model}}}}}\right)
  \end{align*}

\end{frame}

\begin{frame}{Positional Encodings}
  \centering
  \includegraphics[width=\paperwidth,height=0.95\paperheight,keepaspectratio]{assets/Lecture-6_Introduction_to_Transformers/pdf_page_038.png}
\end{frame}

\begin{frame}{Positional Encodings}
  \centering
  \includegraphics[width=\paperwidth,height=0.95\paperheight,keepaspectratio]{assets/Lecture-6_Introduction_to_Transformers/pdf_page_039.png}
\end{frame}

\begin{frame}{Position Encodings}
  \begin{columns}
    \begin{column}{0.6\textwidth}
      \textbf{Requirements for Positional Encodings???}
    \end{column}
    \begin{column}{0.4\textwidth}
      \centering
      % Placeholder for positional encoding/input embedding diagram
      \includegraphics[width=0.5\paperwidth,height=0.5\paperheight,keepaspectratio]{assets/Lecture-6_Introduction_to_Transformers/pdf_page_040.png}
    \end{column}
  \end{columns}
\end{frame}

\begin{frame}{Position Encodings}
  \begin{columns}
    \begin{column}{0.6\textwidth}
      \textbf{Requirements for Positional Encodings}
      \begin{itemize}
        \item Some representation of time? (like \texttt{seq2seq}?)
        \item Should be unique for each position -- not cyclic
      \end{itemize} 
    \end{column}
    \begin{column}{0.4\textwidth}
      \centering
      % Placeholder for positional encoding/input embedding diagram
      \includegraphics[width=0.5\paperwidth,height=0.5\paperheight,keepaspectratio]{assets/Lecture-6_Introduction_to_Transformers/pdf_page_040.png}
    \end{column}
  \end{columns}
\end{frame}

\begin{frame}{Position Encodings}
  \begin{columns}
    \begin{column}{0.6\textwidth}
      \textbf{Requirements for Positional Encodings}
      \begin{itemize}
        \item Some representation of time? (like \texttt{seq2seq}?)
        \item Should be unique for each position -- not cyclic
      \end{itemize}
      \vspace{0.6em}
      \textbf{Possible Candidates:}
      \begin{align*}
        P_{t+1} &= P_t + \Delta c \\
        P_{t+1} &= e^{t \Delta c} \\
        P_{t+1} &= P_t e^{\Delta c}
      \end{align*}
 
    \end{column}
    \begin{column}{0.4\textwidth}
      \centering
      % Placeholder for positional encoding/input embedding diagram
      \includegraphics[width=0.5\paperwidth,height=0.5\paperheight,keepaspectratio]{assets/Lecture-6_Introduction_to_Transformers/pdf_page_040.png}
    \end{column}
  \end{columns}
\end{frame}

\begin{frame}{Positional Encodings}
  \centering
  \includegraphics[width=\paperwidth,height=0.95\paperheight,keepaspectratio]{assets/Lecture-6_Introduction_to_Transformers/pdf_page_043.png}
\end{frame}

\begin{frame}{Positional Encodings}
  \centering
  \includegraphics[width=\paperwidth,height=0.95\paperheight,keepaspectratio]{assets/Lecture-6_Introduction_to_Transformers/pdf_page_044.png}
\end{frame}

\begin{frame}{Position Encodings}
  \begin{columns}
    \begin{column}{0.6\textwidth}
      \textbf{Requirements for Positional Encodings}
      \begin{itemize}
        \item Some representation of time? (like \texttt{seq2seq}?)
        \item Should be unique for each position -- not cyclic
        \item \textbf{Bounded}
      \end{itemize}

      \textbf{Possible Candidates}
      \[
        P(t + t') = M^{t'} \times P(t)
      \]

    \end{column}
    \begin{column}{0.4\textwidth}
      \centering
      % Placeholder for positional encoding/input embedding diagram
      \includegraphics[width=0.5\paperwidth,height=0.5\paperheight,keepaspectratio]{assets/Lecture-6_Introduction_to_Transformers/pdf_page_040.png}
    \end{column}
  \end{columns}
\end{frame}

\begin{frame}{Position Encodings}
  \begin{columns}
    \begin{column}{0.6\textwidth}
      \textbf{Requirements for Positional Encodings}
      \begin{itemize}
        \item Some representation of time? (like \texttt{seq2seq}?)
        \item Should be unique for each position -- not cyclic
        \item \textbf{Bounded}
      \end{itemize}

      \textbf{Possible Candidates}
      \[
        P(t + t') = M^{t'} \times P(t)
      \]

      \vspace{0.8em}
      {\color{blue}\textbf{M?}}
      \begin{enumerate}
        \item Should be a unitary matrix
        \item Magnitudes of eigen value should be 1 $\,\rightarrow\,$ norm preserving
        \item The matrix can be learnt
        \item Produces unique rotated embeddings each time
      \end{enumerate}
 

    \end{column}
    \begin{column}{0.4\textwidth}
      \centering
      % Placeholder for positional encoding/input embedding diagram
      \includegraphics[width=0.5\paperwidth,height=0.5\paperheight,keepaspectratio]{assets/Lecture-6_Introduction_to_Transformers/pdf_page_040.png}
    \end{column}
  \end{columns}
\end{frame}

\begin{frame}{Positional Encodings}
  \centering
  \includegraphics[width=\paperwidth,height=0.95\paperheight,keepaspectratio]{assets/Lecture-6_Introduction_to_Transformers/pdf_page_045.png}
\end{frame}

\begin{frame}{Positional Encodings}
  \centering
  \includegraphics[width=\paperwidth,height=0.95\paperheight,keepaspectratio]{assets/Lecture-6_Introduction_to_Transformers/pdf_page_048.png}
\end{frame}


% \begin{frame}{Position Encodings: Sine and Cosine Functions}
%   \begin{columns}
%     \begin{column}{0.56\textwidth}
%       \textbf{Requirements for Position Encodings}
%       \begin{itemize}
%         \item Some representation of time? (like seq2seq)
%         \item Should be unique for each position
%         \item Bounded
%       \end{itemize}

%       \vspace{1em}
%       \textbf{Actual Candidates}
%       \begin{itemize}
%         \item \textit{sin}(g(t))
%         \item \textit{cos}(g(t))
%       \end{itemize}
      
%       \vspace{1em}
%       \textbf{Requirements for $g(t)$}
%       \begin{itemize}
%         \item Must have same dimensions as input embeddings
%         \item Must produce overall unique encodings
%       \end{itemize}
%     \end{column}
%     \begin{column}{0.42\textwidth}
%       % Placeholder for illustration, plot, or diagram.
%       \centering
%       \begin{tikzpicture}[scale=0.85]
%         \begin{axis}[
%           width=5.7cm,
%           height=3.8cm,
%           grid=both,
%           axis lines=middle,
%           xlabel={$x$},
%           ylabel={},
%           xmin=0,xmax=2*pi,
%           ymin=-1.2,ymax=1.2,
%           xtick={0,1.5708,3.1416,4.7124,6.2832},
%           xticklabels={$0$, $\frac{\pi}{2}$, $\pi$, $\frac{3\pi}{2}$, $2\pi$},
%           ytick={-1,0,1},
%           yticklabels={$-1$,$0$,$1$},
%           legend style={draw=none, fill=none, at={(0.95,0.85)}, anchor=north east, font=\footnotesize},
%         ]
%           \addplot[domain=0:2*pi, samples=100, smooth, thick, blue] {sin(deg(x))};
%           \addlegendentry{$\sin(x)$}
%           \addplot[domain=0:2*pi, samples=100, smooth, thick, red, dashed] {cos(deg(x))};
%           \addlegendentry{$\cos(x)$}
%         \end{axis}
%       \end{tikzpicture}
 
%     \end{column}
%   \end{columns}
  
%   \vspace{1em}
%   % Placeholder for transformer diagram and sine/cosine plot from slide
%   \begin{minipage}{0.55\textwidth}
%     % Text and bullet points (already provided)
%   \end{minipage}%
%   \begin{minipage}{0.4\textwidth}
%     \centering
%     \fbox{\parbox{0.95\linewidth}{\centering \textit{[Transformer diagram and sine/cosine plot here]}}}
%   \end{minipage}
% \end{frame}

\begin{frame}{Position Encoding}
  \begin{columns}
    \begin{column}{0.6\textwidth}
      For each position, an embedded input is moved the same distance but at a different angle. \textbf{Inputs that are close to each other in the sequence have similar perturbations, but inputs that are far apart are perturbed in different directions.}

      \[
        PE_{(pos,\,2i)} = \sin\left( pos / 10000^{2i/d_{\text{model}}} \right)
      \]
      \[
        PE_{(pos,\,2i+1)} = \cos\left( pos / 10000^{2i/d_{\text{model}}} \right)
      \]

      \vspace{1em}
      \begin{description}
        \item[pos] idx of the token in input sentence
        \item[i] $i^\text{th}$ dimension out of $d$
        \item[$d_\text{model}$] embedding dimension of each token
      \end{description}
      Different calculations for odd and even embedding indices.
 
 

    \end{column}
    \begin{column}{0.4\textwidth}
      \centering
      % Placeholder for positional encoding/input embedding diagram
      \includegraphics[width=0.5\paperwidth,height=0.5\paperheight,keepaspectratio]{assets/Lecture-6_Introduction_to_Transformers/pdf_page_040.png}
    \end{column}
  \end{columns}
\end{frame}

\begin{frame}{Position Encoding}
  \centering
  \includegraphics[width=0.8\paperwidth,height=0.8\paperheight,keepaspectratio]{assets/Lecture-6_Introduction_to_Transformers/pdf_page_050.png}
\end{frame}

\begin{frame}{Position Encoding}
  \centering
  \includegraphics[width=0.8\paperwidth,height=0.8\paperheight,keepaspectratio]{assets/Lecture-6_Introduction_to_Transformers/pdf_page_051.png}
\end{frame}

\begin{frame}{Alternatives to Positional Encoding}
  \begin{itemize}
    \item[\textcolor{red}{$\blacktriangleright$}] \textbf{Learnable position embeddings:}
    \begin{itemize}
      \item Position embeddings are parameters learned during training.
      \item Allow the model to adapt position representations to the task.
    \end{itemize}
    \vspace{0.5em}
    \item[\textcolor{red}{$\blacktriangleright$}] \textbf{Rotary Position Embeddings (RoPE):}
    \begin{itemize}
      \item Encodes positions by rotating query and key vectors in multi-head attention.
      \item Enables better extrapolation to longer sequences.
    \end{itemize}
    \vspace{0.5em}
    \item[\textcolor{red}{$\blacktriangleright$}] \textbf{ALiBi (Attention with Linear Biases):}
    \begin{itemize}
      \item Adds linear biases to attention scores based on distance between tokens.
      \item Encourages attention to nearby tokens without explicit position embeddings.
    \end{itemize}
  \end{itemize}
\end{frame}

